%%=============================================================================
%% Conclusie
%%=============================================================================

\chapter{Conclusie}%
\label{ch:conclusie}

In deze bachelorproef werd onderzocht welke factoren bepalend zijn voor wedstrijdstatistieken in de Jupiler Pro League,
en in welke mate data-gedreven machine-learningmodellen en menselijke voetbalexperts overeenstemmen in hun voorspellingen en verklaringen.
In dit afsluitende hoofdstuk worden de onderzoeksvragen beantwoord op basis van de gevonden resultaten,
wordt de wetenschappelijke bijdrage van het werk toegelicht, en worden enkele aanbevelingen voor toekomstig onderzoek geformuleerd.

\section{Antwoord op de onderzoeksvragen}

\subsection*{Deelvraag 1: Welke team- en wedstrijdgerelateerde factoren zijn geassocieerd met het wedstrijdresultaat?}

Uit de historische data-analyse blijkt dat recente vorm en rangverschil de sterkste statistisch significante
associaties vertonen met het wedstrijdresultaat. Het thuisvoordeel is eveneens duidelijk aanwezig: 43,28\% van
de wedstrijden wordt gewonnen door de thuisploeg, tegenover 32,37\% door de bezoeker. Bij een rangverschil
van tien plaatsen of meer wint de hoger gerangschikte ploeg in 78,9\% van de gevallen. Onderlinge confrontaties (H2H)
tonen wel een statistisch significante associatie, maar de praktische voorspellende waarde blijft beperkt (38,29\% nauwkeurigheid als enkel criterium).

\subsection*{Deelvraag 2: Welke factoren beïnvloeden het verloop van wedstrijdstatistieken zoals doelpunten, schoten en kaarten?}

Voor het aantal doelpunten is de speelstijlcombinatie aanvallend versus verdedigend de enige factor
met een statistisch significant effect (p = 0,0139). Wedstrijden in deze configuratie leveren gemiddeld 3,19 doelpunten op,
tegenover een algemeen gemiddelde van 2,86. Het model voor doelpunten bevestigt dit en voegt rangverschil
en vormverschil toe als bijkomende verklarende features.

\vspace{5pt}

Voor het aantal schoten kon geen enkele individueel getoetste factor een statistisch significante invloed aantonen.
Het model slaagt er evenmin in schoten betrouwbaar te voorspellen (R² = $-$0,09), wat aangeeft dat het spelverloop
op dit vlak sterk afhankelijk is van de wedstrijdcontext.

\vspace{5pt}

Voor het aantal kaarten bleek enkel het type wedstrijd als degradatieduel een statistisch significant effect te hebben (p = 0,0399),
met gemiddeld 5,12 kaarten tegenover 4,13 in het algemeen. Dit wijst erop dat extreme motivationele omstandigheden
een meetbare invloed hebben op de spelintensiteit.

\subsection*{Deelvraag 3: Spelen externe contextuele factoren zoals weersomstandigheden, stadioncapaciteit en tijdstip een rol?}

De externe contextuele factoren hebben over het algemeen een beperkte invloed op de wedstrijdstatistieken.
Noch temperatuur noch neerslag vertonen een statistisch significant effect op het aantal schoten, doelpunten of kaarten in de uitgevoerde t-tests. 
Stadioncapaciteit vertoont daarentegen een statistisch significant, 
maar zeer zwak positief verband met het aantal overtredingen van het uitteam (r = 0,079, p < 0,001), 
waarbij de effectgrootte beperkt is. 
Dit verband kan worden geïnterpreteerd als een mogelijke benadering voor verhoogde competitieve druk in grotere stadions met een intensiever thuispubliek, 
al laat de correlatie geen causale interpretatie toe. Het tijdstip van de wedstrijd (avond versus dag, weekend versus weekdag)
en weervariabelen spelen een beperkte rol in de modellen, maar zijn geen bepalende voorspellers.

\subsection*{Deelvraag 4: Hoe kunnen de belangrijkste voorspellende factoren geïdentificeerd en geïnterpreteerd worden?}

Via feature importance-analyses en SHAP-waarden konden de meest bepalende variabelen per doelstatistiek
worden geïdentificeerd. Voor het wedstrijdresultaat zijn dit vorm- en rangverschil. Voor doelpunten
domineren opnieuw rangverschil, maar ook temperatuur (mogelijks met nog andere achterliggende factoren die hier gepaard mee gaan). 
Voor schoten is de speelstijlcombinatie de belangrijkste groep features. Voor kaarten neemt eveneens de speelstijlcombinatie de eerste plaats in.
De SHAP-analyse toont daarnaast de richting van deze effecten: zo gaat een groter rangverschil gepaard
met een hoger voorspeld aantal doelpunten, terwijl kleinere stadions samenhangen met meer voorspelde schoten.

\subsection*{Deelvraag 5: In welke mate kunnen interpreteerbare technieken zoals feature importance inzicht bieden?}

Feature importance en SHAP-waarden bleken bruikbaar om modelverklaringen te vergelijken met menselijke redeneringen.
De combinatie van beide technieken geeft zowel het relatieve belang van variabelen (feature importance)
als de richting en grootte van het effect (SHAP) inzichtelijk weer.
Dit maakt het mogelijk om te beoordelen of een model plausibele verklaringen geeft, of dat het vooral schijnverbanden oppikt,
zoals bij temperatuur als voor seizoenseffecten.

\subsection*{Deelvraag 6: In welke mate komen data-gedreven verklaringen overeen met menselijke expertinschattingen?}

Er is een gedeeltelijke maar geen volledige overeenstemming. Op het vlak van het wedstrijdresultaat
zijn model en experts grotendeels eens: recente vorm en rangverschil worden door beide als belangrijkste
determinanten gezien. Ook de speelstijlcombinatie wordt zowel door het model als door experts als relevante
factor erkend voor doelpunten en kaarten.

\vspace{5pt}

Tegelijk zijn er duidelijke verschillen. Experts overschatten de voorspellende waarde van H2H-resultaten,
een bias die vermoedelijk samenhangt met de cognitieve beschikbaarheidsheuristiek bij opvallende matchen.
Omgekeerd leunt het model sterk op temperatuur als verklarende variabele, terwijl dit voor experts weinig relevant is.
Bij schoten en kaarten overschatten experts ook de invloed van derby's en rivaliteitswedstrijden,
iets wat de data niet ondersteunen. Degradatiewedstrijden zijn als enige factor statistisch aantoonbaar
voor kaarten, maar worden noch door experts noch door het model als dominant beschouwd.

\section{Algemene conclusie}

Het centrale antwoord op de onderzoeksvraag luidt als volgt: \textit{recente vorm, rangverschil en speelstijl zijn de meest bepalende factoren voor wedstrijdstatistieken in de Jupiler Pro League. Data-gedreven modellen en menselijke experts komen op grote lijnen overeen, vooral voor het voorspellen van wedstrijdresultaten, maar verschillen duidelijk van elkaar bij factoren zoals H2H, temperatuur en derby’s.}

\vspace{5pt}

De onvoorspelbaarheid van voetbal blijft een fundamentele uitdaging.
De beste accuraatheid die bereikt werd op het wedstrijdresultaat bedraagt 55,97\%, wat in lijn ligt
met wat gerapporteerd wordt in de internationale literatuur en aangeeft dat het voorspelbaar
plafond in voetbal samenvalt met zijn onzekere, menselijke aard.

\vspace{5pt}

De bijdrage van dit onderzoek ligt in de expliciete vergelijking van drie analyseperspectieven
op eenzelfde dataset van de Jupiler Pro League. Door data, modellen en expertkennis naast elkaar te plaatsen,
wordt niet alleen de beperkte maar reële voorspellende kracht van interpreteerbare modellen gedemonstreerd,
maar worden ook verschillen tussen menselijke inschattingen en historische data zichtbaar gemaakt.
Dit heeft praktische relevantie voor clubs, analisten en journalisten die data-gedreven inzichten willen integreren
in hun besluitvorming, zonder de beperkingen ervan uit het oog te verliezen.

\section{Aanbevelingen voor toekomstig onderzoek}

Op basis van de bevindingen en beperkingen van dit onderzoek kunnen een aantal richtingen voor toekomstig onderzoek worden geformuleerd.

\vspace{5pt}

\textbf{Opname van scheidsrechtersdata.} Het ontbreken van informatie over de aanwezige scheidsrechter
is vermoedelijk het grootste gebrek in het kaartenmodel. Toekomstig onderzoek dat de persoonlijke
stijl en strengheid van scheidsrechters als variabele opneemt, zou de verklaringskracht op dit vlak
aanzienlijk kunnen verbeteren.

\vspace{5pt}

\textbf{Spelersniveau en teamsamenstelling.} Het toevoegen van individuele spelersdata, blessures, transfers
en opstellingen zou de modellen aanzienlijk kunnen verrijken. Deze data zijn weliswaar kostbaar of moeilijk
automatisch te verzamelen, maar zijn beschikbaar via commerciële databronnen.

\vspace{5pt}

\textbf{Grotere expertsteekproef.} Met slechts zes experts zijn de enquêteresultaten indicatief maar niet representatief.
Een toekomstige studie met een grotere en meer diverse groep, inclusief scouts, coaches en data-analisten
bij Belgische profclubs, zou een robuuster beeld geven van hoe menselijke expertise zich verhoudt tot algoritmische verklaringen.

\vspace{5pt}

\textbf{Andere competities en transferleerbaarheid.} Het is interessant om na te gaan in hoeverre de gevonden
verbanden en modelprestaties overdraagbaar zijn naar andere competities met een ander competitieniveau of Europese topcompetities.
Ook werd in dit onderzoek geen rekening gehouden met eventuele bijkomende nationale en internationale wedstrijden, die de druk en vermoeidheid beïnvloeden.

\vspace{5pt}

\textbf{Complexere modelarchitecturen.} Terwijl dit onderzoek bewust koos voor interpreteerbare modellen,
zou het gebruik van ensemble-stacking of gespecialiseerde neurale netwerken kunnen nagaan
of hogere complexiteit ook tot betere voorspellingen leidt, en hoe die resultaten via post-hoc methoden
zoals SHAP inzichtelijk gemaakt kunnen worden.